\documentclass[a4paper]{book}

\usepackage[T1]{fontenc}
\usepackage[french]{babel}
%\usepackage{ProfLycee}
% Définitions des variables
\newcommand{\class}{1NSI}
\newcommand{\datee}{28 janvier 2025}

\usepackage{../../../preambule/preambule_new}


% *** Réglage des headers ***
%\newcommand{\classe}{1NSI}
%\newcommand{\date}{28 janvier 2025}
%\fancyhead[L]{1NSI --- Lycée Victor Hugo}
%\fancyhead[R]{28 janvier 2025}
% *** Réglage des headers ***



%%%%%%%%%%%%%%%%%%%%%% A RANGER %%%%%%%%%%%%%%%%%%%%%%ù


\pagestyle{DS_FP}
\fancyhead[C]{Devoir 6}
   %%%%%%%%%%%%%%%%%%%%%%%%%%%%%%%%%%%%%%%%%%%%%%%%%%%%%%%%
\begin{document}

\emph{Ce devoir porte sur le début du chapitre ARSE11 ainsi que sur les algorithmes de base python sur les listes et dictionnaires}

%\begin{center}
%  {\scshape\LARGE Devoir surveillé \No6\par}
%  \vspace{0.5cm}
%\end{center}

\NomPrenomNote{}

\exodss{1}
\begin{enumerate}
\minipag{0.55}{
\item  Compléter la modélisation de l'architecture de Von Neumann suivante avec les mots:
  \medskip
  \fbox{ \texttt{Mémoires}} \qquad  \fbox{\texttt{Unité arithmétique et Logique}}
 
 \fbox{\texttt{Périphériques}} \qquad   \fbox{\texttt{Unité de contrôle }}
	
	\medskip
	
\item Quelles sont les caractéristiques des mémoires volatiles. Donner un exemple:

\boxguess{4}{1}
}% Fin minipage
\hfill
\minipag{0.42}{
\pictures{1}{1NSI-ARSE11-Architecture_Ordinateurs.jpg}
}

\end{enumerate}

\exodss{2}

On considère la liste suivante:


\begin{CodePythontexAlt}[Largeur=0.95\linewidth, Centre, 
TaillePolice = \normalsize, 
EspacementVertical = 1.2,
Lignes=false,
]{flush right}
liste = [6, 9, 8, 5]
\end{CodePythontexAlt}


\end{document}
	





\medskip
\begin{exerciceinterro}{\hspace{2cm}\dotfill/\fbox{2 points}}{}


\begin{CodePythontexAlt}[Largeur=17cm,
TaillePolice = \normalsize, 
EspacementVertical = 1.2,
Lignes=false]{}
liste = [6, 9, 8, 5]
\end{CodePythontexAlt}


\begin{enumerate}

\item Donner la commande permettant d'insérer la valeur \verb+7+ à la suite de la liste de manière à obtenir:

\begin{CodePythontexAlt}[Largeur=17cm,
TaillePolice = \normalsize, 
EspacementVertical = 1.2,
Lignes=false]{}
[6, 9, 8, 5, 7]
\end{CodePythontexAlt}

	\zonereponse{1cm}

\item Donner la commande permettant de modifier la valeur 8 de manière à la remplacer par \verb+-1+ et obtenir 

\begin{CodePythontexAlt}[Largeur=17cm,
TaillePolice = \normalsize, 
EspacementVertical = 1.2,
Lignes=false]{}
[6, 9, -1, 5, 7]
\end{CodePythontexAlt}

	\zonereponse{1cm}


\end{enumerate}
\end{exerciceinterro}

\medskip

\begin{exerciceinterro}{\hspace{2cm}\dotfill/\fbox{2 points}}{}

On considère le dictionnaire suivant:

\begin{CodePythontexAlt}[Largeur=17cm,
TaillePolice = \normalsize, 
EspacementVertical = 1.2,
Lignes=false]{}
dico = {'Alice' : 7, 'Bob' : 9}
\end{CodePythontexAlt}


\begin{enumerate}
\item Donner la commande permettant d'afficher la valeur correspondante à la clé \verb+'Bob'+.

	\zonereponse{1cm}

\item Donner la commande permettant d'ajouter la clé \verb+'Charlie'+ avec la valeur \verb+4+ et obtenir 

\begin{CodePythontexAlt}[Largeur=17cm,
TaillePolice = \normalsize, 
EspacementVertical = 1.2,
Lignes=false]{}
assert dico == {'Alice' : 7, 'Bob' : 9, 'Charlie' : 4}
\end{CodePythontexAlt}

	\zonereponse{1cm}

\end{enumerate}
\end{exerciceinterro}

\newpage
\pagestyle{DS_LP}

\begin{exerciceinterro}{\hspace{2cm}\dotfill/\fbox{1 points}}{}


 Compléter la fonction \verb+renvoie_liste_valeurs(dico)+ qui prend un dictionnaire en paramètre et renvoie une liste contenant toutes les valeurs du dictionnaire. On utilisera obligatoirement une boucle.


\begin{CodePythontexAlt}[Largeur=17cm,
TaillePolice = \normalsize, 
EspacementVertical = 1.2,
Lignes=false]{}
def renvoie_liste_valeurs(dico):
	''' dico (dict)  dictionnaire {str : int}
	sortie (list) une liste de valeurs entières (int) '''
	
	
	
	
		
		
			

dico = {'Alice':7, 'Bob':9, 'Charlie':7, 'David':10}	
assert renvoie_liste_valeurs(dico) == [7, 9, 7, 10]
	
\end{CodePythontexAlt}

\end{exerciceinterro}

\begin{exerciceinterro}{\hspace{2cm}\dotfill/\fbox{3 points}}{}

\begin{enumerate}
	\item Compléter la fonction \verb+est_positifs(L)+ qui prend une liste de valeurs entières et renvoie \verb+True+ si toutes ces valeurs sont positives ou nulles et \verb+False+ sinon.
	
\begin{CodePythontexAlt}[Largeur=17cm,
TaillePolice = \normalsize, 
EspacementVertical = 1.2,
Lignes=false]{}
def est_positifs(L):
    
      
    
    
    
    
	
assert est_positifs([6, 8, -3, 9]) == False
assert est_positifs([6, 8, 3, 9]) == True
\end{CodePythontexAlt}
	

\item Compléter la fonction \verb+renvoie_cle_val_positives(dico)+ qui renvoie une liste contenant les clés de \verb+dico+ dont les valeurs sont toutes positives. On utilisera la fonction \verb+est_positifs(L)+.

\begin{CodePythontexAlt}[Largeur=17cm,
TaillePolice = \normalsize, 
EspacementVertical = 1.2,
Lignes=false]{}
def renvoie_cle_val_positives(dico):
	''' Renvoie une liste contenant les clés de dico dont 
	les valeurs sont des listes d'entiers tous positifs ou nuls '''
	
	
	
	
	
	
		
	
	
	
dico1 = {'Alice' : [5, 8, 6], 'Bob' : [6, -5, 8], 'Charlie' : [9, 4]}	
assert renvoie_cle_val_positives(dico1) == ['Alice', 'Charlie']
\end{CodePythontexAlt}
\end{enumerate}

\end{exerciceinterro}
\end{document}
